% =====================================================================
\chapter{Méthodologie}
\label{chap-methodologie}
% =====================================================================

Ce chapitre décrit le système construit pour mettre à l'épreuve la thèse formulée
au chapitre précédent, et en détaille la formalisation mathématique. Il commence
par établir ce qu'un navigateur expose réellement de la géométrie d'une page
rendue (\cref{sec-rendu_dom}) représentant le préalable technique dont dépendent tous les
choix qui suivent, puis suit l'ordre de la chaîne de traitement~: une page est
rendue une fois par un navigateur sans interface graphique et sa géométrie est
mesurée (\cref{sec-rendu})~; cette géométrie pilote un tuilage adaptatif, dont
la \cref{sec-tuilage} formalise les règles et le modèle de coût
(\cref{sec-tuilage})~; les tuiles sont organisées en un graphe à deux niveaux,
éventuellement étendu à plusieurs pages liées (\cref{sec-graphe})~; un lecteur
vision-langage est entraîné par contraste, hors ligne, à évaluer la pertinence
d'une paire (question, tuile) (\cref{sec-contrastif})~; enfin, à la requête, un
contrôleur en ligne parcourt le graphe sous budget pour décider quels nœuds le
lecteur ouvre réellement (\cref{sec-controleur}). La \cref{sec-implementation}
détaille les choix d'ingénierie qui rendent l'ensemble exécutable, et la
\cref{sec-protocole} la construction des jeux de données évalués au
\cref{chap-resultats}.

\begin{figure}[H]
\centering
\begin{tikzpicture}[
  node distance=0.62cm,
  every node/.style={font=\small},
  stage/.style={draw, rounded corners=2pt, minimum width=7.2cm, minimum height=0.86cm,
                align=center, fill=codebg, draw=rulegray, line width=0.7pt},
  online/.style={stage, fill=accent!6, draw=accent, line width=0.9pt},
  offline/.style={stage},
  arr/.style={-{Stealth[length=2.2mm]}, thick, draw=black!55},
  tag/.style={font=\scriptsize\itshape, text=black!55, align=left}
]
\node[offline] (render) {Rendu de page et extraction géométrique\\[-2pt]
  {\scriptsize Playwright, retrait du bruit d'interface, boîtes englobantes}};
\node[offline, below=of render] (tiling) {Tuilage adaptatif et découpe\\[-2pt]
  {\scriptsize appariement des titres, fusion, chevauchements}};
\node[offline, below=of tiling] (graph) {Graphe multigranulaire\\[-2pt]
  {\scriptsize 1 nœud page $+$ $N$ nœuds élément, 7 relations, corpus multi-pages}};
\node[offline, below=of graph] (train) {Lecteur entraîné par contraste\\[-2pt]
  {\scriptsize QA synthétique, négatifs difficiles, LoRA sur ViT $+$ décodeur}};
\node[online, below=of train] (ctrl) {Contrôleur de preuves budgété\\[-2pt]
  {\scriptsize \anglais{best-first} sous budget, machine à quatre états}};
\node[online, below=of ctrl] (answer) {Synthèse de la réponse\\[-2pt]
  {\scriptsize lecture des tuiles ouvertes, puis génération}};

\draw[arr] (render) -- (tiling);
\draw[arr] (tiling) -- (graph);
\draw[arr] (graph) -- (train);
\draw[arr] (train) -- (ctrl);
\draw[arr] (ctrl) -- (answer);

\node[tag, right=0.3cm of render.east, anchor=west] {\cref{sec-rendu}};
\node[tag, right=0.3cm of tiling.east, anchor=west] {\cref{sec-tuilage}};
\node[tag, right=0.3cm of graph.east, anchor=west] {\cref{sec-graphe}};
\node[tag, right=0.3cm of train.east, anchor=west] {\cref{sec-contrastif}};
\node[tag, right=0.3cm of ctrl.east, anchor=west] {\cref{sec-controleur}};

\node[rotate=90, anchor=south, font=\scriptsize\bfseries, text=black!55]
  at ([xshift=-0.55cm]$(render.west)!0.5!(train.west)$)
  {hors ligne, une fois par page};
\node[rotate=90, anchor=south, font=\scriptsize\bfseries, text=accent]
  at ([xshift=-0.55cm]$(ctrl.west)!0.5!(answer.west)$)
  {en ligne, par requête};
\end{tikzpicture}
\caption[Vue d'ensemble de la chaîne de traitement]{Vue d'ensemble de la chaîne
DOM-Graph RAG}
\label{fig-pipeline}
\end{figure}

% =====================================================================
\section{Préliminaires géométriques : le DOM et le moteur de rendu}
\label{sec-rendu_dom}
% =====================================================================

L'hypothèse centrale de ce travail est que la segmentation nécessaire au RAG
visuel est déjà présente dans le DOM d'une page rendue. Cette section établit
quelle information le navigateur expose, où se situent les pièges, et pose la
notation reprise par le reste du chapitre.

\subsection{Du balisage à la mise en page, et le modèle de boîtes}
\label{sec-rendu_boites}

Le traitement d'une page par un navigateur enchaîne quatre étapes. L'analyse du
HTML produit le DOM, un arbre $T = (V, E)$ dont chaque nœud $v \in V$
correspond à un élément ou à un nœud texte. La cascade attribue ensuite à chaque
élément ses valeurs de propriétés calculées. La mise en forme (\anglais{reflow})
détermine la position et les dimensions exactes de chaque boîte dans un repère
absolu. La peinture produit enfin les pixels.

La mise en forme se termine intégralement avant que le premier pixel ne soit
peint, si bien que la géométrie de tout élément est déjà connue au moment de la
peinture. Un système qui recadre des régions
dans l'image peinte travaille donc en aval d'une source de vérité géométrique
qu'il n'a jamais consultée.

La fonction $\mathcal{B} : V \to \mathbb{R}_{\geq 0}^4$ associe à un élément
rendu sa boîte $\mathcal{B}(v) = (x_v, y_v, w_v, h_v)$ dans le repère absolu du
document, où $(x_v, y_v)$ est le coin supérieur gauche et $(w_v, h_v)$ la
largeur et la hauteur. C'est elle qu'exploite le tuilage de la
\cref{sec-tuilage}, aucune partie de la chaîne ne redérivant une géométrie que
le navigateur a déjà calculée.

Une boîte est faite de quatre régions concentriques, contenu,
\anglais{padding}, bordure et marge. Les API de géométrie rapportent par défaut
celle qui correspond visuellement à l'élément, et $\mathcal{B}$ désigne
implicitement cette région.

Le type d'affichage détermine le comportement de $\mathcal{B}$. Un élément de
type \textbf{bloc} (\code{<p>}, \code{<div>}, \code{<h1>}--\code{<h6>},
\code{<ul>}) a par défaut \code{width: auto}, donc $w_v = w_{\op{parent}(v)}$
quel que soit son contenu. Un titre de trois mots dans une colonne de
\SI{1200}{\pixel} a ainsi une boîte de \SI{1200}{\pixel} de large. Un élément de
type \textbf{en ligne} (\code{<span>}, \code{<a>}, \code{<em>}) ne forme pas une
boîte unique mais autant de rectangles que de fragments de ligne. Un
\code{<table>}, enfin, s'ajuste à son contenu au lieu de remplir son conteneur.
La difficulté décrite ci-dessous ne frappe donc que les blocs de texte.

\subsection{Éléments flottants : l'écart entre boîte de bloc et lignes réelles}
\label{sec-eda_dom_float}

Le cas déterminant pour la mesure des tuiles est celui des éléments flottants.
Un élément déclaré \code{float: right} est retiré du flux normal et poussé sur
le bord droit de son conteneur, et le contenu en ligne des blocs qui le suivent
s'enroule autour de lui. L'effet est asymétrique. Pour un bloc $v$ qui suit un
flottant, $\mathcal{B}(v)$ n'est pas réduite et continue de s'étendre sur
$w_{\op{parent}(v)}$ en passant sous le flottant, alors que seules les lignes de
texte à l'intérieur de $v$ sont raccourcies. La \cref{fig-eda_float} met les
deux géométries en regard.

\begin{figure}[H]
\centering
\begin{subfigure}[b]{0.50\textwidth}
\centering
\begin{tikzpicture}[scale=0.88, every node/.style={font=\scriptsize}]
\draw[draw=rulegray, line width=0.8pt, fill=codebg] (0,0) rectangle (7.0,4.9);
\node[anchor=north west, text=black!55] at (0.08,4.84) {conteneur};
\draw[draw=black!45, fill=black!12, line width=0.7pt] (4.75,2.30) rectangle (6.85,4.50);
\node[align=center, text=black!60] at (5.80,3.40) {élément\\flottant};
\draw[draw=accent, dashed, line width=1.0pt] (0.18,0.52) rectangle (6.85,4.16);
\foreach \y in {3.94,3.60,3.26,2.92,2.58}{
  \draw[draw=codestr, line width=0.8pt, fill=codestr!12] (0.32,\y-0.12) rectangle (4.55,\y+0.12);
}
\foreach \y in {2.10,1.76,1.42,1.08}{
  \draw[draw=codestr, line width=0.8pt, fill=codestr!12] (0.32,\y-0.12) rectangle (6.68,\y+0.12);
}
\draw[draw=codestr, line width=0.8pt, fill=codestr!12] (0.32,0.62) rectangle (2.60,0.86);
\draw[draw=accent, dashed, line width=1.0pt] (0.20,-0.65) -- (0.95,-0.65);
\node[anchor=west, text=black!65] at (1.05,-0.65) {\code{getBoundingClientRect}};
\draw[draw=codestr, line width=0.8pt, fill=codestr!12] (0.20,-1.32) rectangle (0.95,-1.10);
\node[anchor=west, text=black!65] at (1.05,-1.21) {\code{getClientRects}, nœuds texte};
\end{tikzpicture}
\caption{Les deux géométries d'un même bloc}
\label{fig-eda_float_schema}
\end{subfigure}
\hfill
\begin{subfigure}[b]{0.455\textwidth}
\centering
\begin{tikzpicture}
\node[anchor=south west, inner sep=0] (img)
  {\includegraphics[height=6.0cm, trim=179 38 180 2, clip]{extracted-000.png}};
\begin{scope}[x={(img.south east)}, y={(img.north west)}]
  \draw[draw=accent, dashed, line width=0.9pt] (0.024,0.102) rectangle (0.982,0.239);
  \draw[draw=codestr, line width=0.9pt] (0.038,0.105) rectangle (0.655,0.236);
\end{scope}
\end{tikzpicture}
\caption{La liste de références d'une page réelle}
\label{fig-eda_float_reel}
\end{subfigure}
\caption[Boîte de bloc et rectangles de lignes]{\textbf{(a)}~Un élément flottant à droite raccourcit les
lignes des blocs qui le suivent sans réduire leur boîte de bloc, laquelle
continue de passer sous lui. \textbf{(b)}~La boîte de bloc de la liste de références, en tirets,
s'étend sous l'infoboîte, alors que l'union de ses rectangles de lignes s'arrête à son bord}
\label{fig-eda_float}
\end{figure}

La correction employée en \cref{sec-rendu_textflow} se formalise ainsi. Soit $v$
un élément de texte coulant et $N_1, \dots, N_k$ ses nœuds texte descendants.
\code{getClientRects}, appliqué à un objet \code{Range} construit sur chaque
$N_i$, retourne les rectangles de ligne effectivement rendus
$\{r_{i,1}, \dots, r_{i,m_i}\}$. Chaque $r$ est un quadruplet
$(x^{(0)}, y^{(0)}, x^{(1)}, y^{(1)})$ de coordonnées de coins. La boîte serrée
est l'enveloppe rectangulaire de leur union.
\begin{equation}
\mathcal{B}_{\text{serré}}(v) = \Big[
  \min_{i,j} x^{(0)}_{i,j},\;
  \min_{i,j} y^{(0)}_{i,j},\;
  \max_{i,j} x^{(1)}_{i,j} - \min_{i,j} x^{(0)}_{i,j},\;
  \max_{i,j} y^{(1)}_{i,j} - \min_{i,j} y^{(0)}_{i,j}
\Big].
\label{eq-boite_serree}
\end{equation}
Le \code{Range} doit être construit sur les nœuds texte et non sur
l'élément lui-même. Un \code{Range} englobant tout $v$ retournerait en plus le
rectangle de sa boîte pleine largeur, ce qui ramènerait l'union à
$\mathcal{B}(v)$. Un nœud texte n'a par construction aucune boîte propre, et ses
rectangles ne peuvent donc être que des lignes réellement rendues.

En l'absence de flottant, $\mathcal{B}_{\text{serré}}(v) = \mathcal{B}(v)$, les
lignes remplissant la largeur du bloc. En présence d'un flottant qui intersecte
les rangées de $v$, la boîte serrée est strictement plus étroite, et l'écart
d'aire
\begin{equation}
\Delta A(v) = \op{aire}\big(\mathcal{B}(v)\big) - \op{aire}\big(\mathcal{B}_{\text{serré}}(v)\big) \;\geq\; 0
\label{eq-aire_perdue}
\end{equation}
mesure le fond blanc qu'un recadrage sur $\mathcal{B}(v)$ soumettrait en pure
perte à l'encodeur visuel. C'est ce terme, sommé sur toutes les balises de texte
coulant d'une page, que la mesure serrée de la \cref{sec-rendu_textflow}
élimine, et dont la \cref{sec-res_a_pixels} rapporte l'effet agrégé sur le
pixels par page.

Pour un \code{<table>}, un \code{<figure>} ou un \code{<img>}, en revanche,
bordure, fond et marge intérieure font partie du contenu visuel. La boîte brute
$\mathcal{B}(v)$ y reste la mesure correcte, et $\Delta A(v) = 0$ par
convention.

\subsection{Les deux API de géométrie et le repère absolu}
\label{sec-eda_dom_api}

Le navigateur expose deux méthodes de mesure. \code{getBoundingClientRect}
retourne $\mathcal{B}(v)$, et \code{getClientRects} la liste des rectangles
effectivement occupés, base de $\mathcal{B}_{\text{serré}}$ ci-dessus. Toutes
deux sont relatives au \anglais{viewport} et non au document, alors que le
recadrage d'une capture pleine page exige le repère absolu. Le décalage de
défilement courant $(s_x, s_y)$ leur est donc ajouté.
\begin{equation}
\mathcal{B}_{\text{abs}}(v) = \mathcal{B}_{\text{viewport}}(v) + (s_x, s_y, 0, 0).
\label{eq-repere_absolu}
\end{equation}
Cette conversion est appliquée systématiquement à l'extraction
(\cref{sec-rendu}). La profondeur d'un nœud dans $T$ fournit par ailleurs un
ordre de traitement commode pour résoudre l'imbrication
(\cref{sec-tuilage_elagage}), un parent étant toujours à une profondeur
strictement inférieure à celle de ses descendants.

\subsection{Ordre du DOM et ordre visuel}
\label{sec-eda_dom_ordre}

L'ordre des nœuds dans $T$ est celui dans lequel l'auteur a écrit le balisage,
et non l'ordre de lecture visuel, le positionnement CSS pouvant faire diverger
les deux arbitrairement. Le déduire du parcours de $T$ produirait, sur une mise
en page non triviale, une séquence qui ne correspond pas à ce qu'un lecteur
humain suit. La relation d'ordre de lecture
(\cref{sec-graphe_ordre}) est donc calculée par une heuristique spatiale sur les
coordonnées $\mathcal{B}_{\text{abs}}$ plutôt que par un parcours de $T$. Le DOM
fournit $\mathcal{B}$ et les types d'éléments, pas un ordre de lecture total.

% =====================================================================
\section{Rendu de page et extraction géométrique}
\label{sec-rendu}
% =====================================================================

\subsection{Rendu sans interface graphique}

La page est chargée sans interface graphique, via Playwright, avec une largeur
de fenêtre d'affichage de \SI{2048}{\pixel} et un facteur d'échelle de
périphérique de \num{1}. La hauteur de la fenêtre, fixée à \SI{1000}{\pixel},
est sans effet, la capture pleine page produisant une image de la hauteur totale
du document rendu, soit plusieurs dizaines de milliers de pixels sur les
articles les plus longs (\cref{sec-tuilage_sections}).

La largeur de \SI{2048}{\pixel} égale la hauteur maximale de tuile
(\cref{sec-tuilage_regles}), de sorte que la plus grande unité produite est un
carré de \num{2048} pixels de côté et qu'un seul paramètre borne l'aire maximale
soumise à l'encodeur. Elle n'améliore en revanche pas la résolution du texte, la
taille en pixels d'un caractère ne dépendant pas de la largeur de fenêtre à
facteur d'échelle constant. Élargir le rendu ajoute de la marge, non de la
définition.

Cette marge n'est pas neutre pour le \cref{chap-resultats}. La colonne de
contenu occupe environ \SI{40}{\percent} de la largeur de rendu
(\cref{fig-tuiles_reelles}), si bien qu'une bande de grille est majoritairement
du fond blanc là où une tuile mesurée serrée (\cref{sec-rendu_textflow}) ne
l'est pas. Le rapport de budget de pixels de la \cref{sec-res_a_pixels} est
mesuré à cette largeur et décroît avec elle.

\paragraph{Ordre des opérations.} Le chargement attend la condition
\code{networkidle}, et l'extraction géométrique est évaluée après la capture
pleine page, la condition d'attente étant réévaluée entre les deux. La capture
fait en effet défiler le document et déclenche le chargement différé des images,
qui modifie la hauteur des blocs qui les contiennent. Une mesure prise avant ce
défilement décrirait une mise en page antérieure au chargement, et les
recadrages seraient décalés d'autant sans qu'aucune métrique ne le signale.

\subsection{Retrait du bruit d'interface}

Avant la capture, un ensemble de sélecteurs CSS retire du document le bruit
d'interface, c'est-à-dire les barres de navigation, pieds de page, panneaux
latéraux du logiciel de wiki, bandeaux d'avertissement, liens d'édition de
section et boîtes de navigation. Ces éléments ne portent aucune information sur
le sujet de la page et produiraient des tuiles coûteuses en pixels et
systématiquement non pertinentes.

{\sloppy
La liste combine des sélecteurs génériques (\code{nav}, \code{footer},
\code{header}, \code{script}, \code{style}) et des sélecteurs propres à MediaWiki
(\code{\#mw-navigation}, \code{.vector-header}, \code{.navbox},
\code{.mw-editsection}, \code{.noprint}). Cette seconde catégorie est une
dépendance explicite au type de site et constitue la principale limite de portée
des résultats quantitatifs (\cref{sec-res_limites}). Elle porte sur le nettoyage
et non sur la logique de tuilage.
\par}

\subsection{Extraction des boîtes englobantes}

L'extraction est un unique appel de fonction JavaScript évalué dans le contexte
de la page. Pour un ensemble configurable de balises (\code{h1} à \code{h4},
\code{p}, \code{table}, \code{figure}, \code{img}, \code{ul}, \code{ol},
\code{blockquote}, \code{pre}), le script relève, pour chaque élément
correspondant, sa position absolue au sens de l'\cref{eq-repere_absolu}, ses
dimensions, un aperçu de son texte, sa profondeur dans l'arbre DOM et, le cas
échéant, les liens hypertextes qu'il contient.

\paragraph{Extraction sélective des liens.} Les liens ne sont relevés que pour
les éléments \code{p}. Une page encyclopédique porte quelques dizaines de liens
dans son texte et plusieurs centaines dans ses tableaux, boîtes de navigation et
listes de références, dont le suivi ferait croître le coût du parcours de corpus
(\cref{sec-graphe_corpus}) sans apporter de contenu pertinent. Les liens
conservés sont résolus en URL absolue via \code{document.baseURI}, les ancres
internes et les schémas non HTTP étant écartés.

\subsection{Mesure serrée des balises de texte coulant}
\label{sec-rendu_textflow}

Recadrer sur $\mathcal{B}$ plutôt que sur $\mathcal{B}_{\text{serré}}$ produit
deux défauts mesurables, l'un et l'autre conséquences de l'asymétrie formalisée
en \cref{sec-eda_dom_float}.

Le premier est un double comptage. Recadrer sur $\mathcal{B}(v)$ capture les
pixels de l'élément flottant voisin en plus du contenu visé, et ces pixels
apparaissent alors dans deux unités de récupération, celle du bloc et celle de
l'infoboîte. Le budget est payé deux fois, et les scores de pertinence de deux
tuiles sont alimentés par le même contenu.

Le second est un coût. Une tuile de texte recadrée sur $\mathcal{B}(v)$ consomme
la même surface qu'une bande de grille quel que soit son contenu réel, soit
\num{2058000} pixels en moyenne contre \num{157170} pour une tuile recadrée sur
$\mathcal{B}_{\text{serré}}(v)$ (\cref{sec-res_a_pixels}). Le cas dominant est
la liste de références en fin d'article, rendue en \code{<ol>}, dont la boîte
traverse toute la largeur du conteneur alors que ses lignes occupent une colonne
étroite (\cref{fig-eda_float_reel}).

Le tuilage recadre donc sur $\mathcal{B}_{\text{serré}}(v)$ pour les balises
dont la boîte de bloc est structurellement plus large que leur contenu
(\code{h1} à \code{h4}, \code{p}, \code{blockquote}, \code{ul}, \code{ol}),
désignées balises de texte coulant et configurables. Un \code{TreeWalker}
parcourt les nœuds texte non vides de $v$, un \code{Range} est construit sur le
contenu de chacun, et l'\cref{eq-boite_serree} est évaluée sur l'union des
rectangles retournés par \code{getClientRects}. La granularité du \code{Range}
est imposée par la contrainte établie en \cref{sec-eda_dom_float}. Les autres
balises conservent $\mathcal{B}(v)$, avec $\Delta A(v) = 0$.

% =====================================================================
\section{Tuilage adaptatif guidé par le DOM}
\label{sec-tuilage}
% =====================================================================

Le tuilage instancie pour le cas web le découpeur
$C : \mathcal{D} \mapsto \mathcal{U}$ de la \cref{sec-eda_rag}. La fonction
$C_{\text{DOM}}$ définie ci-dessous, construite à partir de $\mathcal{B}$ et
$\mathcal{B}_{\text{serré}}$ (\cref{sec-rendu_dom}), remplace la grille de
hauteur fixe de PixelRAG. Elle transforme une liste d'éléments géolocalisés $E$
et une capture pleine page $I$ en une liste de tuiles
$T = C_{\text{DOM}}(E, I)$, chacune formée d'un recadrage d'image accompagné de
ses coordonnées, de la balise DOM d'origine, d'un aperçu textuel, des
identifiants des éléments sources et des liens agrégés.

Le nombre de tuiles retenues par unité de page est borné à \num{25}. Cette
borne est héritée du coût de génération des paires question-réponse
(\cref{sec-protocole}) et restreint la portée de l'évaluation~A
(\cref{sec-res_limites_plafond}). Trois traitements préparatoires précèdent le
tuilage proprement dit.

\subsection{Élagage des éléments imbriqués}
\label{sec-tuilage_elagage}

Les balises retenues s'imbriquent, un \code{<img>} dans un \code{<figure>}, un
\code{<p>} dans une cellule de \code{<table>}, une liste \code{<ul>} dans un
\code{<blockquote>}. Sans traitement, chaque niveau produirait sa propre tuile
et le contenu de l'enfant apparaîtrait deux fois.

L'élagage traite les éléments par profondeur DOM croissante et ne compare
chaque candidat qu'aux éléments déjà retenus. Le candidat est écarté si sa
boîte est entièrement contenue dans celle d'un élément retenu, à une tolérance
d'un pixel près qui absorbe les arrondis de rendu. L'ordre par profondeur
croissante garantit qu'un parent est examiné avant ses descendants, donc qu'un
enfant n'est jamais comparé à un autre enfant lui-même susceptible d'être
élagué.

\subsection{Résolution des chevauchements partiels}
\label{sec-tuilage_chevauchement}

L'élagage ne traite que les inclusions complètes. Deux éléments indépendants
peuvent se chevaucher sans que l'un contienne l'autre, même après la mesure
serrée de la \cref{sec-rendu_textflow}. Une infoboîte flottante et le bloc qui
partage sa rangée s'intersectent ainsi sur quelques lignes, là où le texte
reprend la pleine largeur sous l'élément flottant. Les deux sont conservés,
mais les tuiler séparément dupliquerait la région commune. Ils sont donc réunis
en un seul recadrage couvrant leur union. Le critère de regroupement est le
rapport de l'aire d'intersection à l'aire du plus petit des deux éléments,

\begin{equation}
\rho(a, b) = \frac{\op{aire}(B_a \cap B_b)}
                  {\min\!\big(\op{aire}(B_a),\, \op{aire}(B_b)\big)},
\label{eq-overlap}
\end{equation}

deux éléments étant regroupés dès que $\rho(a,b) \geq \num{0.15}$. Le seuil
strictement positif absorbe les recouvrements de quelques pixels dus aux marges
et aux bordures.

Le regroupement s'effectue par composantes connexes sur cette relation, au
moyen d'une structure \anglais{union-find}. La transitivité est nécessaire. Si
$a$ chevauche $b$ et $b$ chevauche $c$ sans que $a$ et $c$ ne se chevauchent,
les trois doivent finir dans le même groupe, faute de quoi la tuile de
$\{a,b\}$ et celle de $\{b,c\}$ partageraient les pixels de $b$.

Les titres sont exclus du regroupement, la règle d'appariement de la
\cref{sec-tuilage_regles} exigeant qu'ils restent des éléments isolés en
attente.

\subsection{Les trois règles de tuilage}
\label{sec-tuilage_regles}

Les groupes sont triés de haut en bas puis de gauche à droite et traités
séquentiellement selon trois règles.

\paragraph{Appariement des titres.} Un titre n'est pas une tuile autonome.
Celle d'un \code{<h2>} seul ne contiendrait que quelques mots sur un fond
blanc, coûteuse en unité de récupération et sémantiquement vide, puisqu'un
titre sans son contenu ne répond à aucune question. Un titre rencontré est donc
mis en attente, et le groupe suivant est fusionné avec lui dans une tuile
unique dont la balise composée (\code{h2+p}) enregistre la fusion, réutilisée
en \cref{sec-graphe_hierarchie}. Deux cas font exception et émettent un titre
seul, deux titres consécutifs et un titre en fin d'unité de page.

\paragraph{Fusion des éléments courts.} Un élément non titre dont la hauteur
est inférieure à \SI{80}{\pixel} est mis en tampon. Le tampon s'accumule
jusqu'à ce que l'étendue verticale de son union franchisse ce minimum, puis est
vidé en une tuile unique de balise \code{merged}. Sans cette règle, une
succession de courtes lignes de catégories en bas de page produirait une
avalanche de micro-tuiles quasi vides, consommant chacune une entrée d'index et
un appel de génération. La tuile~7 de la \cref{fig-tuiles_reelles} en donne le
résultat. Un groupe issu d'un chevauchement
(\cref{sec-tuilage_chevauchement}) échappe à cette règle et n'est jamais mis en
tampon, la non-duplication de pixels primant sur l'heuristique de fusion.

\paragraph{Découpe des éléments démesurés.} Un élément dont la hauteur dépasse
\SI{2048}{\pixel}, un très long tableau ou une liste de références étendue, est
tranché en $\lceil h / h_{\max} \rceil$ sous-tuiles de hauteur égale. C'est le
seul endroit où le tuilage recourt à une découpe indifférente au contenu, à
l'intérieur d'un unique élément DOM devenu trop grand pour l'encodeur. Cette
découpe est une version strictement plus étroite du mode de défaillance de la
grille aveugle, signalée comme telle en \cref{sec-res_limites}.

L'\cref{alg-tuilage} donne la procédure complète.

\begin{algorithm}[t]
\caption{Tuilage adaptatif guidé par le DOM}
\label{alg-tuilage}
\DontPrintSemicolon
\KwData{éléments $E$ avec boîtes englobantes, capture pleine page $I$,
        hauteurs $h_{\max}$ et $h_{\min}$, seuil $\rho_{\min}$}
\KwResult{liste de tuiles $T$}
$E \leftarrow \proc{ÉlaguerImbriqués}(E)$ \tcp*{profondeur DOM croissante}
$G \leftarrow \proc{GrouperChevauchements}(E, \rho_{\min})$ \tcp*{\anglais{union-find}}
trier $G$ par $(y, x)$ du premier membre\;
$T \leftarrow \varnothing$~; $\textit{tampon} \leftarrow \varnothing$~; $\textit{titre} \leftarrow \bot$\;
\ForEach{$g \in G$}{
  \If{$|g| = 1$ \textbf{et} $g_1$ est un titre}{
    \textsc{ViderTampon}()~; \textsc{ViderTitre}()\;
    $\textit{titre} \leftarrow g_1$
    \textbf{continuer}\;
  }
  \If{$\textit{titre} \neq \bot$}{
    $T \leftarrow T \cup \proc{RecadrerDécouper}(\{\textit{titre}\} \cup g)$\;
    $\textit{titre} \leftarrow \bot$~; \textbf{continuer}\;
  }
  \If{$|g| > 1$}{
    \textsc{ViderTampon}()\;
    $T \leftarrow T \cup \proc{RecadrerDécouper}(g)$
    \textbf{continuer}\;
  }
  \If{$\op{hauteur}(g_1) < h_{\min}$}{
    $\textit{tampon} \leftarrow \textit{tampon} \cup \{g_1\}$\;
    \If{$\op{étendue}(\textit{tampon}) \geq h_{\min}$}{\textsc{ViderTampon}()}
  }
  \Else{
    \textsc{ViderTampon}()\;
    $T \leftarrow T \cup \proc{RecadrerDécouper}(\{g_1\})$\;
  }
}
\textsc{ViderTampon}()~; \textsc{ViderTitre}()\;
\Return{$T$}\;
\end{algorithm}

\noindent
\textsc{RecadrerDécouper}$(S)$ recadre $I$ sur l'union des boîtes de $S$, puis
lui applique la découpe ci-dessus si la hauteur obtenue dépasse $h_{\max}$.
\textsc{ViderTampon} et \textsc{ViderTitre} émettent respectivement une tuile
fusionnée pour le tampon accumulé et, dans les deux cas d'exception ci-dessus,
une tuile pour un titre resté sans contenu suivant.

\paragraph{Complexité et comptages.} L'élagage compare chaque candidat aux
éléments déjà retenus et le regroupement examine toutes les paires, ce qui
place les deux étapes en $O(n^2)$ dans le pire cas, où $n$ est le nombre
d'éléments extraits. La page la plus dense du corpus en compte \num{278} avant
élagage. La chaîne des comptages médians par unité de page s'établit à
9 éléments extraits, 7 après
élagage, 7 groupes, pour 6 tuiles. Ces deux
étapes restent négligeables devant le rendu de page et, à plus forte raison,
devant les appels au modèle vision-langage (\cref{sec-res_couts}). Une
structure d'indexation spatiale serait l'optimisation naturelle sur des
documents d'un ordre de grandeur plus denses.

\subsection{Modèle de coût : tokenisation visuelle et effet du tuilage}
\label{sec-eda_cout_visuel}
 
Le tuilage guidé par le DOM produit davantage d'unités que la grille aveugle
(\cref{sec-res_a_pixels}), et rien ne garantit d'avance que cela coûte moins
cher à l'encodeur qui les consomme. Le modèle de coût qui suit tranche la
question.
 
Un modèle vision-langage à tour de vision de type \anglais{Vision
Transformer}~\citep{vit2021} découpe une image en \anglais{patches} carrés de
côté fixe $p$ et les traite comme une séquence de tokens par
auto-attention~\citep{transformer2017}. Pour une image de dimensions
$w \times h$, le nombre de tokens visuels produits est, à un facteur de
regroupement près,
\begin{equation}
n_{\text{tokens}}(w, h) \;\approx\; \frac{w \cdot h}{p^2} \;=\; \frac{\op{aire}(\mathcal{B})}{p^2},
\label{eq-tokens}
\end{equation}
soit une quantité proportionnelle à l'aire de la boîte recadrée et indifférente
à sa forme. Un recadrage deux fois plus étroit à hauteur égale coûte donc deux
fois moins de tokens. Chaque pixel de $\Delta A(v)$ épargné par la mesure serrée
est ainsi un token visuel que l'encodeur n'a jamais à calculer, sans perte
d'information puisque ce pixel était du fond blanc.
 
Cette proportionnalité n'est pas illimitée, le processeur d'image bornant
l'aire encodée entre un plancher $A_{\min}$ et un plafond $A_{\max}$. Le coût
réel suit donc une fonction d'écrêtage.
\begin{equation}
\begin{split}
    n_{\text{tokens}}^{\text{écrêté}}(\mathcal{B}) = \frac{\op{clamp}\big(\op{aire}(\mathcal{B});\, A_{\min},\, A_{\max}\big)}{p^2},\\
    \quad
    \op{clamp}(a; a_{\min}, a_{\max}) = \min\big(\max(a, a_{\min}), a_{\max}\big)
\end{split}
\label{eq-tokens_clamp}
\end{equation}
Comparer deux stratégies de découpage sur leur aire totale de pixels ne revient
à les comparer sur leur coût en tokens qu'à la condition que leurs unités
respectent $A_{\min} \leq \op{aire}(\mathcal{B}) \leq A_{\max}$. C'est une
condition à vérifier plutôt qu'une propriété acquise.
 
Pour une page découpée en tuiles $T = \{t_1, \dots, t_M\}$, le coût total
d'encodage vaut
$\mathcal{N}(T) = \sum_{i=1}^{M} n_{\text{tokens}}^{\text{écrêté}}(t_i)$. Deux
propriétés de cette somme rendent un découpage plus fin favorable sur cette
famille de modèles.
 
Tant qu'aucune unité n'atteint $A_{\max}$, $\mathcal{N}$ ne dépend que de
l'aire totale occupée, et non du nombre d'unités qui la composent. Produire
davantage d'unités plus petites ne coûte alors pas plus cher tant que l'aire
totale utile ne croît pas.
 
Tant qu'aucune unité ne descend sous $A_{\min}$, aucune ne paie de coût
plancher. Cette seconde propriété tient à la résolution dynamique de
Qwen2-VL~\citep{qwen2vl2024}, employé ici. Sous la résolution canonique fixe
des générations antérieures, chaque unité payait le même coût quelle que soit
sa taille réelle, si bien que le mécanisme exploité ici y serait inopérant.
 
Pour Qwen2-VL, le côté effectif de \anglais{patch} après regroupement vaut
\SI{28}{\pixel}, et le plafond a été porté à
$A_{\max} = \SI{12845056}{\pixel}$, douze fois le défaut de la bibliothèque
(\cref{annexe-hyperparams}). La plus grande unité que la chaîne puisse produire
est une tuile carrée de \num{2048} pixels de côté, soit \SI{4194304}{\pixel},
et la bande de grille en fait \SI{2097152}{\pixel}. Aucune n'atteint le
plafond, donc aucune n'est écrêtée et la condition ci-dessus est vérifiée. Le
\cref{tab-cout_tokens} applique les \cref{eq-tokens,eq-tokens_clamp} aux unités
des deux conditions.
 
\begin{table}[H]
\centering
\small
\caption[Aire et tokens visuels par unité]{Aire et nombre de tokens visuels des
unités produites par les deux conditions de l'évaluation~A}
\label{tab-cout_tokens}
\begin{tabular}{@{}l r r r@{}}
\toprule
 & & \multicolumn{2}{c}{Tokens visuels} \\
\cmidrule(l){3-4}
Unité & Aire (px) & sans écrêtage & sous \SI{1.00}{\mega\pixel} \\
\midrule
Tuile DOM, moyenne & \num{157170} & \num{200} & \num{200} \\
Bande de grille, $2048 \times 1024$ & \num{2097152} & \num{2675} & \num{1280} \\
Tuile DOM maximale, $2048^2$ & \num{4194304} & \num{5350} & \num{1280} \\
\midrule
Par page, tuilage DOM (\num{24.0} unités) & \num{3.79}\,M & \num{4810} & \num{4810} \\
Par page, grille (\num{4.9} unités) & \num{10.17}\,M & \num{13100} & \num{6270} \\
\midrule
\textbf{Rapport grille / DOM} & \textbf{\num{2.7}} & \textbf{\num{2.7}} & \textbf{\num{1.3}} \\
\bottomrule
\end{tabular}
\end{table}
 
Le facteur \num{2.7} de la \cref{sec-res_a_pixels} est donc un rapport de
tokens autant que de pixels. La dernière colonne montre ce qu'il deviendrait au
plafond par défaut, où la bande de grille serait écrêtée quand la tuile DOM ne
le serait pas, ce qui ramènerait le gain à \num{1.3}. L'amplitude du gain dépend donc d'un paramètre de
configuration, ce dont une reproduction doit tenir compte.

\subsection{Découpage en sections des pages démesurées}
\label{sec-tuilage_sections}

Une capture pleine page peut atteindre plusieurs dizaines de milliers de pixels
de hauteur, \num{46000} pour l'article consacré aux Beatles, dont la
discographie et la liste de récompenses dominent la page. À cette échelle, la
génération séquentielle de questions dépasse plusieurs heures par page, sans
point de reprise intermédiaire.

Exclure ces pages serait contraire à l'objectif du projet, qui vise les
documents longs. Elles sont donc découpées en amont du tuilage. Une page dont
la hauteur dépasse dix fois la hauteur maximale de tuile, soit
\SI{20480}{\pixel}, est scindée aux frontières de titres \code{h1} et
\code{h2}. Le contenu précédant le premier titre forme une section
\og préambule\fg{} sans titre. Une section restant trop haute à elle seule, un
\code{h2} suivi d'un tableau de discographie géant, est retranchée par un
budget de hauteur fixe de \SI{20480}{\pixel}. C'est le compromis de
la \cref{sec-tuilage_regles} appliqué un cran au-dessus, entre éléments plutôt
qu'à l'intérieur d'un seul.

Chaque section est ensuite traitée par le reste de la chaîne comme une page à
part entière, avec son propre nœud page et sa propre unité de travail
sauvegardée indépendamment. Le seuil est franchi par \num{12} des \num{19}
articles du corpus, qui produisent \num{185} unités de page
(\cref{tab-res_dataset}). Le découpage est donc le régime courant sur ce
corpus, non un filet de sécurité. L'évaluation~A le désactive et rend chaque
article comme une unité unique, faute de quoi les conditions comparées ne
porteraient pas sur la même étendue de document (\cref{sec-protocole_grille}).

\begin{figure}[H]
\centering
\includegraphics[width=\textwidth]{tiling_preview.png}
\caption[Tuiles et ordre de lecture sur une page réelle]{Tuiles produites sur
\code{en.wikipedia.org/wiki/George\_P.\_Gunn}, colorées par catégorie de contenu
et numérotées dans l'ordre de lecture calculé}
\label{fig-tuiles_reelles}
\end{figure}

% =====================================================================
\section{Construction du graphe multigranulaire}
\label{sec-graphe}
% =====================================================================
 
Les tuiles sont organisées en un graphe orienté multi-arêtes $G = (V, E)$, avec
$V = V_{\text{page}} \cup V_{\text{élém}}$ un nœud \code{page} et un nœud
\code{element} par tuile, et $E = \bigcup_{r \in \mathcal{R}} E_r$ l'union de
sept relations typées, $E_r \subseteq V \times V$. Cinq d'entre elles,
$\mathcal{R}_{\text{MAGE}} \subset \mathcal{R}$, reprennent le graphe de
MAGE-RAG~\citep{magerag2026}. Les deux autres,
$\{\code{links\_to}, \code{continues}\} = \mathcal{R} \setminus \mathcal{R}_{\text{MAGE}}$,
répondent à un besoin propre au contexte web (\cref{sec-graphe_interdoc}).
Le \cref{tab-graphe_relations} les récapitule.
{\sloppy
Le contrôleur ne propage que le long de trois d'entre
elles, $\mathcal{R}_{\text{parcourue}} = \{\code{reading\_order},
\code{links\_to}, \code{continues}\}$, ce que discutent la
\cref{sec-res_synthese} et la \cref{sec-res_limites}.
\par}

\begin{table}[t]
\centering
\small
\caption[Les sept relations du graphe]{Les sept relations du graphe
multigranulaire et leur densité mesurée \textit{(la dernière colonne donne la médiane des arêtes par unité de page)}}
\label{tab-graphe_relations}
\begin{tabularx}{\textwidth}{@{}l l X r@{}}
\toprule
\textbf{Relation} & \textbf{Domaine} & \textbf{Définition} & \textbf{Arêtes} \\
\midrule
\code{contains} & page $\to$ élém. & rattachement plat de chaque tuile à sa page & 6 \\
\code{reading\_order} & élém. $\to$ élém. & successeur dans l'ordre de lecture spatial & 5 \\
\code{layout\_adjacency} & élém. $\leftrightarrow$ élém. & voisinage 2D (gauche, droite, dessus, dessous) & 8\\
\code{section\_hierarchy} & titre $\to$ contenu & arbre des sections respectant l'imbrication \code{h1}--\code{h4} & 4\\
\code{semantic\_neighbor} & élém. $\leftrightarrow$ élém. & proximité lexicale TF-IDF au-delà d'un seuil \textit{(désactivée par défaut)} & \num{0} \\
\midrule
\code{links\_to} & élém. $\to$ page & lien hypertexte vers un autre document & 8 \\
\code{continues} & page $\to$ page & section suivante d'une page découpée & 1 \\
\bottomrule
\end{tabularx}
\end{table}

\subsection{Ordre de lecture spatial}
\label{sec-graphe_ordre}
 
L'ordre de lecture n'est pas l'ordre du DOM, pour la raison établie en
\cref{sec-eda_dom_ordre}. Il est calculé par une heuristique spatiale sur les
coordonnées mesurées, dans la lignée des travaux d'analyse de mise en page qui
traitent l'ordre de lecture comme une prédiction plutôt que comme une lecture de
l'arbre~\citep{layoutreader2021}. Les tuiles sont d'abord regroupées en rangées.
Partant de la tuile la plus haute non encore assignée, toutes les tuiles dont
l'ordonnée est à moins de \SI{20}{\pixel} de la sienne forment une rangée.
Chaque rangée est ensuite triée de gauche à droite, et les rangées sont
concaténées de haut en bas.
 
La tolérance de \SI{20}{\pixel} est le seul paramètre libre de l'heuristique.
Elle absorbe les décalages de ligne de base entre un titre et un élément
flottant voisin sans fusionner deux rangées distinctes.

\subsection{Voisinage de mise en page}
\label{sec-graphe_adjacence}
 
L'ordre de lecture aplatit la structure 2D en une chaîne, ce qui rend deux
tuiles côte à côte indiscernables de deux tuiles empilées. La relation
\code{layout\_adjacency} restitue ce voisinage en réutilisant le même
regroupement en rangées. Elle produit deux types d'arêtes.
 
\begin{itemize}
\item Les voisins gauche et droite, entre tuiles consécutives d'une même rangée.
\item Les voisins dessus et dessous, entre une tuile et la tuile de la rangée
suivante qui la recouvre horizontalement. En cas de plusieurs candidates, celle
dont le centre horizontal est le plus proche est retenue.
\end{itemize}
 
Chaque arête est insérée dans les deux sens, car la règle du centre le plus
proche appliquée en sens inverse ne désignerait pas nécessairement la même
voisine. La symétrie est une décision de construction et pas une propriété de la
règle.
 
Le second type capture les voisinages que l'ordre linéaire fait diverger. Le bas
d'une colonne de gauche et le haut de la colonne de droite sont consécutifs dans
l'ordre de lecture mais visuellement éloignés, et inversement.

\subsection{Arbre des sections}
\label{sec-graphe_hierarchie}

La relation \code{section\_hierarchy} construit un arbre de sections, distinct du
rattachement plat de \code{contains}. Un parcours de l'ordre de lecture maintient
une pile de titres ouverts, chacun associé à son niveau. À la rencontre d'un titre
de niveau $\ell$, la pile est dépilée jusqu'à ce que son sommet ait un niveau
strictement inférieur à $\ell$~; une arête est créée depuis ce sommet vers le
nouveau titre, qui est empilé. Une tuile de contenu est rattachée au sommet
courant. Un \code{h3} est ainsi rattaché au dernier \code{h2} ouvert, non au
dernier \code{h1}, et le contenu précédant le premier titre ne produit aucune
arête.

L'appariement de la \cref{sec-tuilage_regles} interagit avec ce parcours. Une tuile
de balise composée \code{h2+p} est traitée comme un titre de niveau~2 --- c'est le
premier composant qui porte le niveau --- et comme le contenu de cette même
section, qui n'engendre donc pas d'arête vers elle-même. En pratique l'arbre est
moins dense qu'avec des titres isolés~: la relation n'apparaît qu'entre titres, et
vers les tuiles de contenu qui suivent une section déjà ouverte sans porter
elles-mêmes de titre. La densité mesurée figure au
\cref{tab-graphe_relations}~; la \cref{sec-res_limites} discute ce que le
contrôleur en tire.

\subsection{Voisinage sémantique}
\label{sec-graphe_semantique}
 
La cinquième relation de MAGE-RAG relie des éléments proches sémantiquement,
indépendamment de leur position. Elle est implémentée par similarité cosinus sur
vecteurs TF-IDF entre les aperçus textuels des tuiles, avec un seuil de
\num{0.2}.
 
Deux garde-fous l'encadrent. Une sélection gloutonne borne le degré afin que les paires
au-delà du seuil soient triées par similarité décroissante, puis retenues tant
qu'aucune des deux tuiles n'a atteint trois voisins sémantiques. Cela évite un
graphe excessivement dense sur une page au contenu répétitif comme une
discographie ou un tableau de résultats sportifs. Cette relation est désactivée par défaut pour servir de condition d'ablation. Cette comparaison n'a pas été conduite ici
(\cref{sec-res_limites}), et toutes les mesures du \cref{chap-resultats} portent
donc sur des graphes à six relations sur sept. La relation inactive est celle
que ce travail hérite de MAGE-RAG. Les deux relations qui lui sont propres sont
actives, et ce sont elles qui produisent l'effet mesuré en
\cref{sec-res_ablation_c}.
 
\subsection{Relations inter-documents}
\label{sec-graphe_interdoc}
 
Les deux relations ajoutées répondent à un besoin que MAGE-RAG n'a pas,
puisqu'il opère sur un document unique déjà délimité.
 
\paragraph{\code{links\_to}.} Cette relation relie un nœud élément à une page
externe, à partir des liens extraits en \cref{sec-rendu}. Un second budget s'y
applique, puisque seules les $k_{\text{liens}} = 3$ premières tuiles porteuses
de liens, dans l'ordre de lecture de l'unité de page, voient leurs liens
conservés. Sur une page non découpée, l'ordre de lecture commence par
l'introduction de l'article, et ce budget retient donc les liens qui définissent
le sujet. Sur une section issue du découpage de la \cref{sec-tuilage_sections},
il retient les premiers liens de cette section, sans garantie de saillance
comparable. Au-delà, les liens sont ignorés pour borner le facteur de
branchement du parcours de corpus.
 
\paragraph{\code{continues}.} Cette relation chaîne les sections d'une page
découpée dans leur ordre d'origine. Elle est sémantiquement distincte de
\code{links\_to}, un lien menant à un autre document quand une continuation mène
à la suite du même. Le contrôleur les traite néanmoins symétriquement, puisque
dans les deux cas franchir l'arête signifie sortir du périmètre de la page
courante.
 
\subsection{Graphe de corpus multi-pages}
\label{sec-graphe_corpus}
 
À partir d'une page d'amorce, le système parcourt jusqu'à $m = 3$ pages
atteignables par ses arêtes \code{links\_to}. Le parcours est d'un seul saut et
non récursif, car suivre les liens de proche en proche ferait croître le nombre
de pages exponentiellement, pour un bénéfice incertain sur des questions qui
portent rarement à plus d'un saut du sujet (\cref{sec-res_limites}).
 
Chaque page est traitée indépendamment par la même chaîne, puis les graphes sont
fusionnés. Deux mécanismes rendent la fusion possible.
 
Le premier est le préfixage par espace de noms. Chaque page utilise
indépendamment les identifiants \code{page}, \code{tile\_0000} et suivants. Tous
les nœuds d'un sous-graphe sont donc préfixés avant fusion par un identifiant
court dérivé de son URL, ce qui garantit l'unicité sans coordination.
 
Le second est le point d'entrée. Le graphe d'une page expose un nœud
page canonique comme cible d'un lien entrant, le nœud unique si la page n'a pas
été découpée et la première section sinon. Lors de la fusion, chaque nœud
\code{external\_page} tenant lieu de substitut est remplacé par le point
d'entrée réel. Ce recâblage transforme une collection de graphes de pages en un
graphe de corpus navigable, et rend possible l'expansion inter-documents du
contrôleur (\cref{sec-controleur_expansion}).
 
\subsection{Rendu structuré pour le lecteur}
\label{sec-graphe_xml}
 
Un sous-graphe en portée pour une requête peut être sérialisé en XML, dans
l'esprit de l'étape de rendu structuré multimodal de MAGE-RAG. Un élément
\code{<page>} porte l'URL et le titre, et contient un élément \code{<element>}
par tuile directement rattachée, avec son identifiant, sa balise DOM, son état
courant et son texte. Ce format donne au modèle lecteur une vue structurée de
ce qui a été retenu et de ce qui a été écarté, plutôt qu'une concaténation de
fragments sans provenance. Il est employé par l'application de
démonstration (\cref{sec-implementation}), tandis que les quatre évaluations du
\cref{chap-resultats} transmettent les preuves sous la forme concaténée de la
\cref{sec-controleur_lecture}.
 
Les identifiants \code{tile\_NNNN} suivent l'ordre de production du tuilage, qui
trie les groupes par position. Il coïncide avec l'ordre de lecture de la
\cref{sec-graphe_ordre} sauf lorsque deux tuiles d'une même rangée visuelle ont
des ordonnées distinctes.
 
\begin{lstlisting}[float=t, language=XML, caption={Rendu XML d'un graphe de page
après passage du contrôleur pour la requête \og When and where was George Purnell
Gunn born?\fg{}},
label={lst-xml}, morekeywords={page,element,url,title,id,tag,state}]
<page url="en.wikipedia.org/wiki/George_P._Gunn" title="George P. Gunn">
  <element id="tile_0000" tag="merged" state="inactive">...</element>
  <element id="tile_0003" tag="p" state="opened">
    Gunn was born on October 11, 1903, in Winona, Mississippi, ...
  </element>
  <element id="tile_0004" tag="h2+p" state="pruned">...</element>
  ...
</page>
\end{lstlisting}

% =====================================================================
\section{Entraînement contrastif du lecteur}
\label{sec-contrastif}
% =====================================================================
 
Un modèle vision-langage générique sait décrire une tuile, mais rien ne le
prépare à juger si elle permet de répondre à une question donnée. Il y est
entraîné par contraste, en quatre temps. Les données sont d'abord construites,
par génération synthétique de paires question-tuile puis minage de négatifs
difficiles. Le modèle génératif est ensuite converti en encodeur, puis adapté
par LoRA. Une perte InfoNCE est enfin optimisée sur les plongements obtenus.
 
\subsection{Génération synthétique de paires question-tuile}
\label{sec-contrastif_qa}
 
Pour chaque tuile retenue, l'image seule est présentée au modèle vision-langage
d'inférence (\code{qwen3-vl}, \cref{annexe-hyperparams}), avec la consigne de
formuler une question factuelle et autonome dont la réponse est visible dans
cette image et uniquement dans cette image. L'invite exige une question
spécifique, portant sur un nombre, un nom, une date, une définition ou une
légende, et une réponse impossible à deviner sans avoir vu la tuile. Si la tuile
ne contient aucune information exploitable, le modèle répond exactement
\code{SKIP}. La réponse attendue est un objet JSON strict, ce qui rend l'analyse
syntaxique déterministe. L'invite complète figure en \cref{annexe-invites}.
 
C'est la transposition à des tuiles d'image de la recette de génération
synthétique de requêtes établie par InPars~\citep{inpars2022} et
Promptagator~\citep{promptagator2022}.
 
Toutes les tuiles retenues sont sauvegardées sur disque, y compris celles ayant
reçu \code{SKIP}, car une telle tuile reste un candidat valide comme négatif
difficile pour une question portant sur une autre tuile de la même page. Sur le
corpus, \num{67} tuiles sur \num{1784} reçoivent \code{SKIP}
(\cref{tab-res_dataset}).
 
\subsection{Minage de négatifs difficiles}
\label{sec-contrastif_mining}
 
Pour chaque question, toutes les autres tuiles de la même page sont classées par
similarité cosinus TF-IDF avec la question, et les $K = 2$ meilleures qui
franchissent un filtre anti-faux-négatif sont retenues. Le scoreur lexical est
choisi pour son coût, puisqu'il ne demande ni GPU ni modèle à charger alors que
le minage s'exécute sur des milliers de tuiles. Sa configuration est reportée en
\cref{annexe-hyperparams}.
 
Le filtre anti-faux-négatif opère en deux temps, de coût très différent.
 
Le premier est lexical et gratuit. Tout candidat dont l'aperçu textuel
contient déjà la réponse verbatim est rejeté, ce qui attrape les faux négatifs
littéraux, par exemple deux blocs mentionnant la même date.
 
Le second, optionnel, est un juge vision-langage. Le filtre lexical ne
détecte que les recouvrements textuels et laisse passer un candidat exprimant la
même information par un synonyme, ou visuellement, comme une couleur dans un
graphique, un logo ou une valeur lue sur un axe. L'image candidate est alors
soumise au modèle avec la question et la réponse attendue. Ce filtre coûte un
appel par candidat non écarté en amont. Il est désactivable, et les jeux de
données du \cref{chap-resultats} sont construits avec le filtre lexical seul.
 
Le minage conserve une question même lorsque moins de $K$ candidats franchissent
le filtre, et l'exigence de $K$ négatifs n'est appliquée qu'au chargement du jeu
de données pour l'entraînement. Sur les \num{1339} exemples construits,
\num{1282} portent les deux négatifs attendus, \num{27} n'en portent qu'un et
\num{30} aucun (\cref{tab-res_dataset}). Ces \num{57} derniers sont écartés au
chargement, après le partitionnement de la \cref{sec-contrastif_split}.
 
Le minage opère au sein d'une même page, les identifiants de tuiles n'étant
uniques qu'à cette échelle. L'étendre au corpus exigerait de reprendre en amont
le préfixage par espace de noms de la \cref{sec-graphe_corpus}. Cette extension
est mise en œuvre pour le chemin HotpotQA (\cref{sec-protocole_multihop}), où
les positifs proviennent de deux pages et où elle est obligatoire.
 
\subsection{Du modèle génératif à l'encodeur}
\label{sec-contrastif_encodeur}
 
Le modèle affiné, \code{Qwen2-VL-7B-Instruct}~\citep{qwen2vl2024}, est
génératif. Il transforme du texte et des images en texte et n'expose aucun
vecteur de représentation, alors qu'une perte contrastive en exige un.
 
La technique PromptEOL~\citep{prompteol2024} lève cet obstacle. L'entrée est
suivie d'une courte instruction fixe, reproduite en \cref{annexe-invites}, et
l'état caché de la dernière couche à la position du dernier token est pris comme
plongement. L'instruction contraint le modèle à comprimer le sens de l'entrée en
une position unique.
 
Le point qui rend la technique adaptée ici est qu'elle est indifférente à la
modalité. Le plongement d'une question textuelle et celui d'une image de tuile
sont calculés de la même manière, avec la même instruction, seule l'entrée qui
la précède changeant. Question et tuile vivent donc par construction dans le
même espace, sans tête de projection additionnelle à apprendre. En notant
$f_\theta$ cette fonction d'encodage, $\mathbf{z}_q = f_\theta(q)$ pour une
question et $\mathbf{z}_u = f_\theta(u)$ pour une tuile $u$, la perte
contrastive de la \cref{sec-contrastif_loss} porte exclusivement sur $\theta$, à
travers les LoRA appliqués aux deux tours (\cref{sec-contrastif_lora}).
 
\paragraph{Où se joue la compréhension visuelle.} La tour de vision de
$f_\theta$ reçoit une séquence de \anglais{patches} (\cref{eq-tokens}). Pour une
tuile issue de la grille aveugle, qui chevauche potentiellement plusieurs
sections hétérogènes (\cref{sec-res_a_attach}), cette séquence mélange les
\anglais{patches} de plusieurs entités sans qu'aucune limite ne les sépare, et
le gradient de l'\cref{eq-infonce} doit pousser $f_\theta$ à extraire un signal
discriminant d'une représentation dont une partie, potentiellement majoritaire,
est hors sujet. Pour une tuile issue du tuilage guidé, chaque \anglais{patch}
appartient par construction à l'unique élément sur lequel porte la question, et
le signal à amplifier occupe toute la séquence. $f_\theta$ reste un Qwen2-VL
générique, sans mécanisme d'attention ni architecture d'encodeur nouvelle. C'est
la frontière de recadrage, déplacée en amont de l'encodeur, qui change la nature
de la difficulté, de juger un mélange de contenus à juger un contenu homogène.
La \cref{sec-res_ablation_b} mesure l'effet de l'entraînement contrastif sur des
tuiles ainsi alignées, mais ce travail ne mesure pas séparément ce que
l'alignement DOM apporte à l'entraînement lui-même, faute d'un entraînement
identique mené sur des tuiles de grille.
 
\paragraph{Chemin d'exécution et modèle affiné.} L'affinage n'emprunte pas le
chemin utilisé pour la génération de questions et la lecture des tuiles
(\cref{sec-contrastif_qa,sec-controleur_lecture}), qui s'adresse à un serveur
Ollama n'exposant qu'une API de génération, sans accès aux poids ni aux
gradients. Il charge le modèle via \code{transformers} et \code{peft}, ce qui
exige un accès complet aux poids et un GPU loué. Le modèle affiné n'est donc pas
celui de l'inférence, et \code{Qwen2-VL-7B-Instruct} a été retenu pour la
stabilité de sa classe \code{peft} au moment de ce travail, quand Qwen2.5-VL et
Qwen3-VL l'étaient moins pour l'affinage par LoRA. C'est un choix d'outillage,
non une revendication d'équivalence entre les deux modèles. Cette proximité de
famille sans identité est reprise en \cref{sec-res_limites_famille} comme risque
de biais à contrôler.
 
\subsection{Adaptation LoRA sur la tour de vision et le décodeur}
\label{sec-contrastif_lora}
 
L'adaptation utilise LoRA~\citep{lora2021} avec un rang $r = 16$, un facteur
d'échelle $\alpha = 32$ et un \anglais{dropout} de \num{0.05}, soit
\num{44 302 336} paramètres entraînables sur les sept milliards du
modèle de base.
 
\paragraph{Cibler les deux tours.} LoRA est appliqué aux projections d'attention
du décodeur (\code{q\_proj}, \code{k\_proj}, \code{v\_proj}, \code{o\_proj}) et à
celles de la tour de vision (\code{qkv}, \code{proj}). Juger si une tuile répond
à une question exige d'affiner simultanément ce qui est vu dans l'image et ce qui
est compris de la question. Geler la tour de vision reviendrait à demander au
décodeur de compenser seul une représentation visuelle inadaptée à la tâche de
récupération. La \cref{sec-eda_contrastif_adaptation} situe ce choix dans les
variantes multimodales de LoRA.
 
\paragraph{Découvrir les modules cibles par introspection.} Les noms de modules
internes de Qwen2-VL varient d'une version de \code{transformers} à l'autre.
Seuls des suffixes sont donc spécifiés, et le code parcourt le modèle chargé pour
découvrir les modules linéaires correspondants. LoRA n'est appliqué qu'aux
modules effectivement trouvés, ce qui rend l'adaptation robuste aux changements
d'architecture interne mais impose de journaliser la liste obtenue à chaque
exécution.
 
\subsection{La perte contrastive}
\label{sec-contrastif_loss}
 
Soit un lot de $B = 4$ questions, chacune accompagnée d'une tuile positive et de
$K = 2$ négatifs difficiles minés. Les plongements des $B$ questions et des
$B(1+K)$ tuiles sont calculés, et la perte minimisée est l'InfoNCE de
l'\cref{eq-infonce} étendue au lot entier,
 
\begin{equation}
\mathcal{L} = -\frac{1}{B}\sum_{i=1}^{B} \log
\frac{\exp\!\big(\mathrm{sim}(\mathbf{z}_{q_i}, \mathbf{z}_{u_i^{+}})/\tau\big)}
     {\displaystyle\sum_{j=1}^{B}\Big[
       \exp\!\big(\mathrm{sim}(\mathbf{z}_{q_i}, \mathbf{z}_{u_j^{+}})/\tau\big)
       + \sum_{l=1}^{K}
       \exp\!\big(\mathrm{sim}(\mathbf{z}_{q_i}, \mathbf{z}_{u_{j,l}^{-}})/\tau\big)
     \Big]},
\label{eq-infonce_lot}
\end{equation}
 
avec une température $\tau = \num{0.05}$. Le dénominateur porte sur les $B(1+K)$
tuiles du lot et non sur les $1+K$ tuiles associées à la question courante, si
bien que chaque positif entre en compétition à la fois avec ses propres négatifs
minés et avec les tuiles de toutes les autres questions du lot, les
\anglais{in-batch negatives} de la \cref{sec-eda_contrastif_infonce}. À $B = 4$
et $K = 2$, le nombre de négatifs effectifs par question,
\begin{equation}
N_{\text{eff}} = B(1+K) - 1 = 11,
\label{eq-negatifs_effectifs}
\end{equation}
passe de \num{2}, les seuls négatifs minés si l'on entraînait avec $B = 1$, à
\num{11}, sans plongement supplémentaire à calculer. Le coût du lot croît
linéairement en $B$, le nombre de comparaisons quadratiquement.
 
\paragraph{Lecture de l'\cref{eq-infonce_lot} comme classification.} Pour une
question $q_i$ fixée, notons $\mathbf{s}_i \in \mathbb{R}^{B(1+K)}$ le vecteur
des similarités $\mathrm{sim}(\mathbf{z}_{q_i}, \cdot)/\tau$ contre les $B(1+K)$
tuiles du lot, et $y_i$ l'indice de la tuile positive $u_i^+$ dans ce vecteur. Le
terme sommé de l'\cref{eq-infonce_lot} est alors exactement l'entropie croisée
catégorielle $-\log \op{softmax}(\mathbf{s}_i)_{y_i}$. Apprendre par contraste
revient donc à résoudre, pour chaque question du lot, un problème de
classification à $B(1+K)$ classes dont une seule est correcte. Le gradient par
rapport à $\mathbf{z}_{q_i}$ se décompose en une force qui rapproche
$\mathbf{z}_{q_i}$ de $\mathbf{z}_{u_i^+}$ et des forces qui l'éloignent de
chaque négatif $u$, chacune pondérée par la probabilité
$\op{softmax}(\mathbf{s}_i)_u$ que le modèle attribue déjà à ce négatif. Un
négatif que le modèle confond encore avec le positif reçoit ainsi un gradient
plus grand qu'un négatif déjà écarté avec confiance.
 
La température $\tau$ contrôle l'acuité de cette pondération. En la faisant
tendre vers $0$, le softmax se concentre sur l'argument maximal et le gradient
se resserre presque entièrement sur le négatif le plus proche du positif, ce qui
accélère l'apprentissage sur les cas difficiles mais amplifie d'autant l'effet
d'un faux négatif introduit par le minage (\cref{sec-contrastif_mining}). À
$\tau = \num{0.05}$, ce compromis penche délibérément du côté de la
discrimination fine, cohérent avec l'objectif de la \cref{sec-contrastif_qa} de
produire des questions dont la réponse n'est visible que dans une tuile précise.
 
\subsection{Contraintes mémoire}
\label{sec-contrastif_memoire}
 
Sur une carte de \SI{80}{\giga\byte}, le chemin de calcul par lots des
plongements a d'abord échoué sur un dépassement de mémoire GPU, avec
\SI{79.16}{\giga\byte} occupés sur les \SI{79.25}{\giga\byte} disponibles.
L'usage est réel et non fragmenté, puisqu'un lot pousse $B(1+K) = 12$ tuiles à
travers le modèle simultanément, soit environ \num{64 000} tokens
visuels au total. Aucune de ces tuiles n'étant écrêtée
(\cref{sec-eda_cout_visuel}), cette longueur croît proportionnellement à l'aire
cumulée du lot.
 
L'activation du \anglais{gradient checkpointing}, qui recalcule les activations
lors de la passe arrière plutôt que de les conserver, résout le problème au prix
de \SIrange{20}{30}{\percent} de temps de calcul supplémentaire. C'est le réglage
par défaut retenu à cette résolution de tuile.
 
\subsection{Partitionnement pour l'évaluation}
\label{sec-contrastif_split}
 
Un partitionnement au niveau des exemples serait trompeur, car deux tuiles
voisines d'une même page produisent des questions fortement corrélées, et un
modèle ayant vu l'une en entraînement bénéficierait indûment sur l'autre en
validation.
 
Le partitionnement est donc au niveau de l'\textbf{article}, en proportion
\num{80}/\num{20} visée puis arrondie à l'article entier, avec des articles
strictement disjoints entre les deux parties. Aucune tuile et aucune question quasi dupliquée ne peut traverser la
frontière. L'adaptateur évalué en \cref{sec-res_ablation_b} est entraîné sur la
seule partie d'entraînement.

% =====================================================================
\section{Contrôleur de preuves en ligne}
\label{sec-controleur}
% =====================================================================
 
À la requête, le contrôleur décide quels nœuds élément le lecteur ouvre
réellement, plutôt que de lui présenter la page entière ou une sélection
top-$k$ figée. Il généralise l'algorithme de contrôle de MAGE-RAG.
 
\subsection{Machine à états}
\label{sec-controleur_etats}
 
Chaque nœud élément porte un attribut d'état, initialisé à \etat{inactive} lors
de la construction du graphe, que le contrôleur fait évoluer à travers la machine
à quatre états de la \cref{fig-etats}. Un nœud \etat{inactive} n'appartient pas
à la frontière courante. Un nœud \etat{active} est candidat à l'ouverture. Un
nœud \etat{opened} a été lu et son contenu fait partie des preuves. Un nœud
\etat{pruned} a été jugé non pertinent, ou laissé de côté faute de budget.
 
\begin{figure}[t]
\centering
\begin{tikzpicture}[
  node distance=1.1cm and 1.9cm,
  st/.style={draw, rounded corners=2pt, minimum width=2.2cm, minimum height=0.75cm,
             align=center, font=\small, draw=rulegray, line width=0.8pt, fill=codebg},
  stf/.style={st, fill=accent!8, draw=accent},
  arr/.style={-{Stealth[length=2.2mm]}, thick, draw=black!60}
]
\node[st] (inactive) {\etat{inactive}};
\node[st, right=of inactive] (active) {\etat{active}};
\node[stf, above right=0.62cm and 2.55cm of active] (opened) {\etat{opened}};
\node[st, below right=0.62cm and 2.55cm of active] (pruned) {\etat{pruned}};
\draw[arr] (inactive) -- node[above, font=\scriptsize]{activation} (active);
\draw[arr] (active) -- node[above, sloped, font=\scriptsize, pos=0.42]{ouverture ($-1$~budget)} (opened);
\draw[arr] (active) -- node[below, sloped, font=\scriptsize, pos=0.42]{élagage (gratuit)} (pruned);
\end{tikzpicture}
\caption[Machine à états du contrôleur]{Machine à états du contrôleur.
L'ouverture d'un nœud consomme une unité de budget et peut activer ses voisins
encore inactifs, le long de \code{reading\_order} sur la même page ou en
franchissant une frontière de document via \code{links\_to} et \code{continues}.
L'élagage est gratuit et s'applique à un nœud actif dont le score est sous le seuil
d'ouverture, puis à tout nœud resté actif lorsque le budget est épuisé. Un nœud
jamais activé reste \etat{inactive}.}
\label{fig-etats}
\end{figure}
 
\subsection{Politique \anglais{best-first} budgétée}
\label{sec-controleur_politique}
 
Le contrôleur cherche un sous-ensemble $O \subseteq V_{\text{élém}}$ de nœuds à
ouvrir sous la contrainte de budget $|O| \leq B$, où $B$ désigne ici le nombre
d'ouvertures autorisées et non la taille de lot de la \cref{sec-contrastif_loss}.
Il maximise informellement la pertinence cumulée $\sum_{n \in O}
\mathrm{score}(n)$ tout en respectant une contrainte d'accessibilité, un nœud ne
pouvant entrer dans $O$ que s'il a d'abord été rendu \etat{active} par
l'ouverture d'un voisin déjà dans $O$, ou par l'amorçage initial.
 
C'est un problème de sélection sous contrainte de cardinalité et de connexité,
pour lequel l'\cref{alg-controleur} implémente une heuristique \anglais{best-first}
gloutonne plutôt qu'une optimisation exacte. Rien ne garantit que $O$ maximise
$\sum \mathrm{score}(n)$ sur les sous-graphes connexes de cardinalité $B$
atteignables depuis l'amorce, seulement que chaque décision locale est
rationnelle au moment où elle est prise. Ce choix reprend celui de MAGE-RAG et
privilégie le calcul en une passe.
 
La politique comporte trois phases. L'amorçage active les
$k_{\text{seed}}$ nœuds inactifs les mieux notés pour la requête, seule phase où
la sélection est globale. La boucle principale prend répétitivement le
nœud actif le mieux noté, l'élague sans consommer de budget si son score est sous
le seuil d'ouverture, et sinon l'ouvre, ce qui consomme une unité de budget et
rend actifs ses voisins encore inactifs qui franchissent le seuil d'activation.
La finalisation élague tout nœud resté actif une fois le budget épuisé,
de sorte qu'aucun nœud ne subsiste dans un état transitoire.
 
\begin{algorithm}[t]
\caption{Contrôle de preuves \anglais{best-first} sous budget}
\label{alg-controleur}
\DontPrintSemicolon
\KwData{graphe $G$, requête $q$, budget $B$, amorce $k_{\text{seed}}$,
        seuils $\theta_{\text{ouv}}$ et $\theta_{\text{act}}$}
\KwResult{liste des nœuds \etat{opened}, en ordre de lecture}
\ForEach{nœud élément $n$ de $G$}{
  $\mathrm{score}(n) \leftarrow \mathrm{sim}\big(q, \operatorname{texte}(n)\big)$ \tcp*{TF-IDF}
}
activer les $k_{\text{seed}}$ nœuds inactifs de plus haut score\;
$b \leftarrow 0$\;
\While{$b < B$ \textbf{et} il existe un nœud \etat{active}}{
  $n \leftarrow \arg\max \mathrm{score}$ parmi les nœuds \etat{active}\;
  \If{$\mathrm{score}(n) < \theta_{\text{ouv}}$}{
    $n \leftarrow \etat{pruned}$ \tcp*{gratuit}
    \textbf{continuer}\;
  }
  $n \leftarrow \etat{opened}$~; $b \leftarrow b + 1$\;
  $N \leftarrow \proc{VoisinsOrdreDeLecture}(n) \cup \proc{VoisinsInterDocuments}(n)$\;
  \ForEach{$n' \in N$ tel que $n'$ est \etat{inactive}}{
    \If{$\mathrm{score}(n') \geq \theta_{\text{act}}$}{$n' \leftarrow \etat{active}$}
  }
}
élaguer tout nœud encore \etat{active}\;
\Return{les nœuds \etat{opened}, triés par ordre de lecture}\;
\end{algorithm}
 
Ce qui distingue cette politique d'une sélection top-$k$ n'est ni le score ni le
budget, identiques, mais le fait que les voisins ne sont considérés qu'après
l'ouverture de leur nœud. La frontière est conditionnée par ce qui s'est déjà
révélé pertinent plutôt que par un classement global figé.
 
\paragraph{Effet de l'égalité des deux seuils.} Les deux seuils valent \num{0.05}
(\cref{annexe-hyperparams}), et leur égalité restreint l'élagage gratuit. Un
nœud activé par un voisin a franchi $\theta_{\text{act}}$, donc son score atteint
aussi $\theta_{\text{ouv}}$ et il ne peut pas être élagué gratuitement. Seules
les amorces, activées sans test de seuil, le peuvent. Le contrôleur n'ouvre donc
moins de $B$ nœuds que pour deux raisons, l'élagage d'une amorce ou l'épuisement
de la frontière, et la \cref{sec-res_c_analyse} mesure \num{4.86} ouvertures en
moyenne pour un budget de \num{5}.
 
\subsection{Expansion inter-documents}
\label{sec-controleur_expansion}
 
Le contrôleur de MAGE-RAG, et la première implémentation à sa suite, ne propagent
l'activation que le long de \code{reading\_order}, relation qui ne franchit
jamais une frontière de page. Un nœud ouvert près de la fin d'une page ne peut
donc activer aucun nœud d'une page liée ni de la section suivante, alors même que
les deux arêtes existent. Sur une question multi-sauts dont les preuves sont
réparties sur deux articles, le second n'est atteignable que si l'amorce globale
y a déjà placé un nœud, c'est-à-dire exactement ce que fait une sélection top-$k$
statique, et moins bien dès que l'amorce en manque un.
 
Ce point est la cause diagnostiquée de l'incapacité du contrôleur à dépasser la
référence statique sur les questions inter-pages (\cref{sec-res_c_avant}). À
l'ouverture d'un nœud, en plus de ses voisins d'ordre de lecture, sont désormais
rendus candidats
 
\begin{itemize}
\item tous les éléments de toute page atteignable depuis ce nœud par une arête
\code{links\_to}, la page cible étant, grâce au recâblage de la
\cref{sec-graphe_corpus}, une page réellement construite et non un substitut,
\item et tous les éléments de toute page atteignable depuis sa page conteneuse
par une arête \code{continues}.
\end{itemize}
 
Dans les deux cas, franchir la porte n'entraîne aucune activation automatique,
chaque élément candidat devant franchir le même seuil $\theta_{\text{act}}$ qu'un
voisin d'ordre de lecture. Le mécanisme est strictement symétrique à celui qui
existait déjà, seul le voisinage considéré s'élargit, et il reste limité à trois
des sept relations du graphe (\cref{sec-graphe}). La \cref{sec-res_ablation_c}
mesure l'effet de ce correctif.
 
\subsection{Score de pertinence et point d'extension}
\label{sec-controleur_scoreur}
 
Par défaut, la pertinence est mesurée par similarité cosinus TF-IDF entre la
requête et le texte de chaque nœud, avec le même scoreur et la même configuration
que le minage de négatifs (\cref{sec-contrastif_mining}). C'est un substitut
destiné à être remplacé.
 
Le contrôleur expose à cet effet un paramètre de scores précalculés. Si un
dictionnaire de scores est fourni, il remplace intégralement le scoreur lexical,
un nœud absent recevant un score nul. Ce point d'extension est conçu pour y
brancher la similarité du lecteur affiné de la \cref{sec-contrastif}, sans
toucher à la politique. Le calcul de ces scores, un passage avant du modèle par
tuile candidate, reste hors de ce module. La \cref{sec-res_limites} explique
pourquoi ce branchement n'a pas été effectué et pourquoi il constitue la
modification la plus prometteuse restant à faire.
 
\subsection{Lecture des tuiles retenues}
\label{sec-controleur_lecture}
 
L'\cref{alg-controleur} retourne une liste de nœuds, non des preuves. Les tuiles
ouvertes sont ensuite relues par le modèle vision-langage, qui en extrait les
informations utiles à la requête. Les preuves finales sont la concaténation, en
ordre de lecture, de ces extractions, et remplacent donc l'aperçu textuel qui a
servi au scorage.
 
La politique a une propriété structurelle que la \cref{sec-impl_concurrence}
exploite. Quel nœud est ouvert ne dépend jamais du contenu lu par le modèle, mais
uniquement des scores figés en amont. La simulation complète de la machine à
états peut donc être déroulée sans aucun appel au modèle, les lectures n'étant
effectuées qu'ensuite, sur une liste déjà connue.
 
% =====================================================================
\section{Implémentation et ingénierie}
\label{sec-implementation}
% =====================================================================
 
\subsection{Architecture modulaire}
 
Le système est écrit en Python, avec un module par étage. Le rendu et l'extraction
reposent sur Playwright, le tuilage sur des recadrages Pillow, le graphe sur
\code{networkx}. Le client vision-langage s'adresse à un serveur Ollama, ce qui
permet de démarrer avec un seul téléchargement de modèle et sans dépendance à
\code{torch}. Le module d'affinage, seul à dépendre de \code{torch},
\code{transformers} et \code{peft}, n'est pas importé par le reste du système et ne
charge ces bibliothèques qu'à l'exécution. La cartographie des modules figure en
\cref{annexe-depot}.
 
Une application web expose la chaîne complète derrière un formulaire
\og URL plus question\fg{}. La requête déclenche l'exécution en direct, puis
affiche la réponse synthétisée, le graphe interactif (catégorie de contenu par
couleur, état du contrôleur par bordure, arêtes filtrables par relation) et les
pages liées, avec un marqueur sur celles portées par une tuile retenue comme
preuve. C'est le seul consommateur du rendu XML de la \cref{sec-graphe_xml}. Chaque
requête recalcule tout, sans corpus préconstruit, et le temps de réponse va de
quelques dizaines de secondes à plusieurs minutes selon la page et le budget.
 
\subsection{Reproductibilité des appels au modèle}
\label{sec-impl_repro}
 
\paragraph{Décodage glouton.} Aucune température n'étant fixée, le serveur
échantillonnait avec ses réglages par défaut, donc à température strictement
positive. Deux appels sur la même image et la même invite pouvaient produire des
réponses différentes, y compris une réponse correcte et une abstention sans aucune
différence dans les tuiles ouvertes. La température est fixée à zéro, et à entrée
fixée la sortie est déterministe.
 
\paragraph{Fenêtre de contexte.} Le modèle d'inférence raisonne avant de répondre,
sur \numrange{1500}{2500} tokens pour une extraction factuelle sur une tuile. Ce
raisonnement s'ajoute à l'invite et aux tokens visuels de la tuile, soit \num{200}
en moyenne et jusqu'à \num{5350} pour une tuile de taille maximale
(\cref{tab-cout_tokens}). Avec la fenêtre de contexte par défaut du serveur, fixée
à \num{4096} tokens, la somme épuisait la fenêtre avant l'émission de la réponse
finale, et le résultat était un contenu vide, silencieusement interprété comme
\code{SKIP} par l'analyseur de réponses. La fenêtre est portée à \num{8192} tokens,
ce qui laisse peu de marge dans le cas le plus défavorable. Cette classe d'erreurs,
une défaillance d'infrastructure qui se déguise en résultat de modèle légitime, est
celle contre laquelle la journalisation des longueurs de séquence protège le mieux.
 
\paragraph{Cache sur disque.} Chaque appel au modèle est mis en cache sur disque,
indexé par un condensat du triplet (image, invite, configuration de décodage). Le
décodage étant glouton, un appel déjà effectué se rejoue sans coût. Ce cache rend
praticable l'itération sur le code d'évaluation, puisque relancer un script après
correction d'un défaut ne repaie pas les heures d'appels déjà effectuées.
 
\subsection{Sauvegarde par unité de travail}
 
L'état est sauvegardé par unité de page, une unité étant une page ou une section
d'une page découpée. Ce grain répond directement au cas de la
\cref{sec-tuilage_sections}, un échec en cours de traitement d'une page démesurée
ne coûtant plus que la section en vol.
 
\subsection{Concurrence}
\label{sec-impl_concurrence}
 
Deux étapes exploitent la concurrence, avec au plus \num{4} appels au modèle
simultanés, et dans les deux cas la légitimité du parallélisme découle d'une
propriété du système.
 
La génération de questions émet ses appels en parallèle au sein d'une page. Ces
appels sont indépendants, chacun ne recevant qu'une image et une invite fixe.
 
La lecture des tuiles ouvertes fait de même, et c'est ici que la propriété de la
\cref{sec-controleur_lecture} devient opérationnelle. Puisque la trajectoire
d'états ne dépend jamais du contenu lu, la liste des tuiles à lire est entièrement
connue avant la première lecture, et ces lectures sont donc mutuellement
indépendantes. Un unique client est réutilisé entre fils d'exécution, chaque appel
étant une requête HTTP indépendante sans état mutable partagé à protéger.
 
\subsection{Tests}
 
La stratégie de test suit une séparation par dépendances, détaillée en
\cref{annexe-depot}. Les tests sans dépendance lourde couvrent la logique
géométrique, la construction du graphe, le minage de négatifs, l'assemblage du jeu
de données, le cache et la concurrence du client, et la compression d'images. Ils
s'exécutent sur des cas synthétiques, sans réseau, sans modèle et sans GPU.
 
Les tests adossés à Playwright vérifient l'extraction DOM de bout en bout contre
une page HTML locale fournie avec le dépôt. Deux cas y sont verrouillés. Un élément
flottant partageant sa rangée avec du texte fixe le comportement de la mesure
serrée de la \cref{sec-rendu_textflow}. Deux éléments en chevauchement partiel,
tous deux conservés, détectent une régression de la \cref{sec-tuilage_elagage} vers
un élagage trop agressif.


%  =====================================================================
\section{Protocole expérimental}
\label{sec-protocole}
% =====================================================================
 
\subsection{Construction du jeu de données synthétique}
 
Un script pilote la chaîne sur une liste de pages, avec reprise sur échec du
téléchargement. Son unité de travail est l'unité de page, et chaque unité est
sauvegardée indépendamment dans un fichier au format JSON lignes. Trois fichiers
en résultent, un manifeste associant chaque identifiant de tuile à son image, les
paires question-réponse générées, et les exemples contrastifs assemblés avec leurs
négatifs. Cette reprise ne couvre pas un échec survenant après la génération
des questions, dont l'\cref{sec-res_limites_dataset} rapporte l'effet.
 
Le plafond de \num{25} tuiles par unité de page (\cref{sec-tuilage}) borne le coût
de génération. Il borne aussi la portée de l'évaluation~A, ce qu'analyse
l'\cref{sec-res_limites_plafond}.
 
\subsection{Reproduction de la référence en grille}
\label{sec-protocole_grille}
 
Pour l'ablation du tuilage (\cref{sec-res_ablation_a}), un module réimplémente la
référence en grille aveugle indépendamment du module de tuilage. La page rendue est
tranchée en bandes de \SI{1024}{\pixel} de haut à la largeur de rendu, ce qui
reproduit la hauteur de tuile documentée par PixelRAG mais non sa largeur, plus
étroite (\cref{sec-eda_pixelrag}). Cet écart de largeur est à porter au compte de
la comparaison de budget de pixels, pour la raison exposée en \cref{sec-rendu}. Une
troisième condition découpe le texte extrait de la même région de page en fenêtres
plates de \num{400} caractères, sans aucune notion de frontière d'élément. C'est la
référence textuelle standard que l'introduction critique et que ni PixelRAG ni
MAGE-RAG ne mesurent directement.
 
Les trois conditions sont construites sur des articles rendus avec un seuil de
découpage en sections relâché par rapport à celui de la
\cref{sec-tuilage_sections}. Cette évaluation retient une hauteur de page maximale
de \num{25} fois la hauteur maximale de tuile plutôt que \num{10}, soit
\SI{51200}{\pixel} contre \SI{20480}{\pixel} pour le reste du corpus. Ce seuil
propre à l'évaluation, indépendant de celui qui protège des coûts de génération le
corpus principal, suffit à ce que chaque article soit traité comme une unité de
page unique, de sorte que les trois conditions portent sur la même étendue de
document. La \cref{sec-res_ablation_a} précise l'exception rencontrée en pratique
sur l'article le plus long du corpus.
 
L'indépendance de ce module vis-à-vis du module de tuilage est un choix
méthodologique, puisqu'elle garantit qu'une modification du tuilage guidé ne peut
pas altérer silencieusement la référence à laquelle on le compare.
 
\subsection{Ablation de résolution}
\label{sec-protocole_resolution}
 
PixelRAG documente la résolution d'image comme levier d'efficacité, sans préciser
ni le filtre de ré-échantillonnage ni les résolutions testées. Un couple pleine
résolution et résolution réduite est donc construit sur les tuiles de ce travail,
par sous-échantillonnage Lanczos à un facteur d'échelle linéaire de \num{0.5}, soit
quatre fois moins de pixels. Les images compressées sont écrites sous les mêmes
noms dans un dossier miroir, de sorte que le manifeste et les exemples contrastifs
restent réutilisables tels quels, seul le dossier racine des images changeant entre
les deux expériences.
 
\subsection{Questions multi-sauts issues de HotpotQA}
\label{sec-protocole_multihop}
 
La génération synthétique de la \cref{sec-contrastif_qa} a un biais structurel,
puisqu'elle demande une question dont la réponse est visible dans une seule tuile.
Un contrôleur omniscient pourrait donc répondre à n'importe laquelle de ces
questions depuis un unique nœud, ce qui ne sollicite jamais le graphe.
 
Cet écart est comblé avec HotpotQA~\citep{hotpotqa2018}, dont chaque question
nomme, via ses phrases de support, deux articles Wikipédia et les phrases exactes
nécessaires dans chacun. Ce signal annoté par des humains est réutilisé plutôt que
régénéré, selon la procédure suivante.
 
\begin{enumerate}
\item Les titres d'articles cités sont convertis en URL Wikipédia et rendus par la
chaîne existante, inchangée.
\item Pour chaque article, les tuiles dont le texte recouvre une phrase de support
sont retenues comme positives. Une même question peut donc avoir des tuiles
positives sur deux pages distinctes.
\item Les négatifs difficiles sont minés sur le vivier réunissant les tuiles des
deux articles, l'extension inter-pages que la \cref{sec-contrastif_mining} laissait
de côté étant ici obligatoire.
\end{enumerate}
 
Une question est abandonnée dans deux cas, plutôt que de produire un exemple
incomplet. Le premier est un article qui ne se rend plus, renommé ou supprimé
depuis le \anglais{dump} Wikipédia de 2017 sur lequel HotpotQA est construit. Le
second est une phrase de support introuvable dans le rendu actuel, l'article ayant
été réécrit.
 
L'attribution d'une phrase de support à une tuile se fait par recouvrement de
tokens filtrés des mots vides, et non par correspondance de sous-chaîne exacte.
Celle-ci survit rarement à une décennie de modifications éditoriales, et l'exiger
ferait chuter le taux de conservation pour une raison sans rapport avec ce qui est
mesuré.